\section{Background}
\label{sec:background}

\subsection{IEEE 754 Subnormal Floating-Point Numbers}
The IEEE 754 standard for single-precision floating-point arithmetic allocates 32 bits into three fields: 1 sign bit ($s$), 8 exponent bits ($e$, with a bias of 127), and 23 significand bits ($f$). For standard normalized values ($e \in [1, 254]$), the representation assumes an implicit leading bit of 1:
\begin{equation}
x_{\text{norm}} = (-1)^s \times (1.f)_2 \times 2^{e - 127}
\end{equation}
The smallest positive normalized single-precision number occurs when $e = 1$ and $f = 0$, yielding $N_{\text{min}} = 1.0_2 \times 2^{-126} \approx 1.175 \times 10^{-38}$.

When the exponent field is all zeros ($e = 0$), the value transitions to a subnormal (or denormal) representation. In subnormal mode, the implicit leading bit shifts from 1 to 0, and the effective exponent is fixed at $-126$:
\begin{equation}
x_{\text{sub}} = (-1)^s \times (0.f)_2 \times 2^{-126}
\end{equation}
The smallest positive subnormal has bit pattern \texttt{0x00000001}, or $2^{-149} \approx 1.4 \times 10^{-45}$. Subnormals support gradual underflow, bridging the gap between $0.0$ and $N_{\text{min}}$.

\subsection{Microarchitectural Timing Asynchrony}
Modern x86 execution units are built for normalized operands, executing in 1 to 4 clock cycles ($\sim 1\text{ ns}$); full hardware support for subnormals across every execution unit would cost significant silicon and power.

Consequently, when an execution pipeline encounters a subnormal operand or generates a subnormal output under standard IEEE 754 compliance (where Flush-to-Zero and Denormals-are-Zero modes are disabled), the fast execution datapath cannot process the value. The CPU flushes the out-of-order execution pipeline and delegates processing to internal microcode assist routines (\texttt{fp\_assists}). Handling gradual underflow via microcode introduces a severe latency penalty of 80 to 200+ cycles ($\sim 30\text{--}70\text{ ns}$) per operation, creating a data-dependent timing side-channel.

\subsection{Pixel Stealing Attack}
Browsers enforce the Same-Origin Policy (SOP): a page may embed a cross-origin image but cannot read its pixels directly, for example through HTML5 canvas extraction. Andrysco et al. ~\cite{andrysco2015subnormal} showed that subnormal timing can bypass the SOP and reconstruct cross-origin pixels:
\begin{enumerate}
    \item \textbf{SVG Filter Processing:} A target cross-origin image is rendered inside an HTML5/SVG container. The attacker applies a custom SVG filter primitive (such as \texttt{feColorMatrix} or \texttt{feConvolveMatrix}) that scales normalized pixel color channel values $I \in [0.0, 1.0]$ by a minute scalar $\alpha \approx 10^{-42}$.
    \item \textbf{Data-Dependent Latency:}
    \begin{itemize}
        \item Non-Black Pixels ($I > 0$): Multiplying a non-zero color intensity by $\alpha$ yields a subnormal result ($\alpha I \approx 10^{-42}$). This repeatedly triggers microcode assist traps during rendering, significantly stalling the browser execution pipeline.
        \item Black Pixels ($I = 0$): Multiplying zero by $\alpha$ yields exact zero ($0.0$), completing immediately on the fast hardware execution path without latency.
    \end{itemize}
    \item \textbf{Pixel Extraction:} By measuring frame render latencies using browser timing mechanisms (such as requestAnimationFrame or CSS repaint callbacks), the attacker classifies individual pixels as black or non-black based on rendering speed. Repeating this process across a target pixel grid progressively leaks private visual content.
\end{enumerate}